In today's fast-paced world, public transportation is rarely used due to its lack of comfort, inefficiency, or failure to meet users' expectations. Common issues include buses not adhering to schedules or difficulties in purchasing tickets. In many places, schedules are neither displayed at stops nor available online. It often happens that buses are overcrowded, and passengers only realize this when they are already at the stop without having alternative travel options. The Bus4U online platform aims to solve these problems.

Bus4U consists of a mobile application, a client website, and an administrator web interface. Its main features include searching for bus routes by specifying the departure point, destination, and time, purchasing tickets online in advance, and viewing departure times broken down by cities and stops. Additionally, bus stops can be viewed on a map filtered by cities and routes, and buses' current positions can be tracked in real time. To ensure these functionalities and manage tickets, a dedicated application runs on POS terminals installed on buses. This app validates tickets and shares live bus information with the central server.

Furthermore, the web application provides an administration interface for bus companies. Here, administrators can manage schedules, stops, and upload routes. Company administrators can also manage employee information and various bus-related data, such as periodic technical inspections and insurance. A statistical interface is also available, allowing company executives to view the number of registered users, sold tickets, and generated revenue in monthly, yearly, or all-time breakdowns.

Future development possibilities include introducing subscriptions and various discounts, as well as tracking bus occupancy in real time. From the perspective of bus companies, monitoring drivers' working hours and implementing automated accounting could also prove beneficial.

This thesis focuses on the software's user interface, including the two websites and the mobile application, while the backend functionality is covered in another thesis.

\textbf{Keywords:} public transport, bus ride, digitalization